\clearpage
\fancyhead[L]{\small PASQUICH. STERNWARTE IN OFEN.}
\hypertarget{native-502-lower}{}\label{native-502-lower}
\begin{center}
\small Göttingische gelehrte Anzeigen. Stück 85. S. 841..843. 28. Mai 1808.\\[3pt]
\rule{.72\linewidth}{.3pt}
\end{center}

Unter dem Titel \emph{Rechenschaft über meine Vorschläge zur Förderung der Astronomie an der königlichen Universitätssternwarte in Ofen} hat der Direktor dieser Sternwarte, Herr \textsc{Joh. Pasquich}, eine kleine Schrift drucken lassen. Sie enthält sehr erfreuliche Nachrichten über das künftige Aufblühen der Astronomie in Ungarn, auf das man hoffen darf.

{\rightskip=0pt plus 3em Die seit 1780 in Ofen bestehende Sternwarte entspricht weder in ihrer Bauart noch mit ihren Instrumenten dem heutigen Stand der Astronomie. Das war umso bedauerlicher, als nach den Statuten der Universität an der Sternwarte außer dem Direktor noch zwei Assistenten und ein Wärter unterhalten werden müssen. Die königliche ungarische Statthalterei, an deren Spitze der aufgeklärte Erzherzog Palatin steht, erkannte, wie wichtig diese Einrichtung für die Wissenschaft und das Land werden könnte, wenn dort tatkräftige und ganz ihrer Wissenschaft ergebene Männer mit besseren Hilfsmitteln ausgestattet würden. Herr Professor \textsc{Pasquich} erhielt daher gleich nach seiner Anstellung als Astronom an dieser Sternwarte von höchster Stelle den Auftrag, Vorschläge zur Anschaffung neuer Instrumente zu machen. Er tat dies nach dem Grundsatz, dass es hier tatsächlich um die ernsthafte Förderung der Wissenschaft gehe, zu der beizutragen sich alle kultivierten Nationen Europas verpflichtet fühlen. Seine Vorschläge wurden sowohl vom Erzherzog Palatin und der königlichen Statthalterei als auch vom Kaiser selbst sogleich ohne jede Einschränkung genehmigt. Zugleich wurde angeordnet, die Instrumente unverzüglich zu bestellen. Er wandte sich daher an den Artilleriehauptmann \textsc{Reichenbach} in München, der dort bekanntlich vor einigen Jahren eine Werkstatt zur Herstellung mathematischer Instrumente eingerichtet hat und nach mehreren Proben den besten englischen Instrumentenbauern ebenbürtig zu sein scheint. Bestellt wurden ein drei Fuß großer Repetitionskreis mit einem dreißig Zoll großen Azimutalkreis und einem silbernen Limbus, ein sechs Fuß langes Meridianfernrohr, eine astronomische Pendeluhr mit Sekundenpendel, eine astronomische Reisependeluhr mit Halbsekundenpendel, ein achtzehn Zoll großer
\SourceLeaf{503}
astronomischer Kreis, ein zwölf Zoll großer Kreis für terrestrische Messungen, ein acht Fuß langes Fernrohr und ein Äquatoreal. Dies alles, einschließlich einiger Reparaturen an anderen Instrumenten, wurde zum mäßigen Preis von 7210 kaiserlichen Gulden oder 8652 Gulden in Reichswährung vereinbart. Außerdem wurde bei Bergrat \textsc{Seyffert} in Dresden noch eine astronomische Pendeluhr mit Sekundenpendel für 360 Taler bestellt. Die Herstellung jener Instrumente wurde zwar durch den Krieg unterbrochen, ist aber jetzt bis auf das acht Fuß lange Fernrohr und das Äquatoreal ganz abgeschlossen. Ein Teil der Instrumente ist schon in Ofen angekommen, die übrigen werden täglich erwartet.\par}

So bleibt jetzt nur noch zu wünschen, dass diese schöne Ausstattung bald einen würdigen Ort findet, an dem sie zum Wohl der Wissenschaft und zur Ehre der ungarischen Nation und ihrer hochgestellten Förderer eingesetzt werden kann. Herr \textsc{Pasquich} hat bereits einen geeigneten Platz für die neue Sternwarte vorgeschlagen. Nur einige besondere Umstände lassen noch ungewiss erscheinen, ob man sich für ihn entscheiden wird. Auf jeden Fall darf man mit Sicherheit erwarten, dass Ungarn bald mit einem Tempel der Urania geschmückt sein wird, zu dem man der ungarischen Nation und der Wissenschaft selbst Glück wünschen kann.
